\section{Introduction}
\label{sec:introduction}

\begin{sloppypar}
When a cache request misses and the cache is full, a system must pick one
resident object to evict before admitting the requested object. Evaluating
whether a given eviction decision was a good one requires knowing what
would have happened to each candidate victim if it, rather than another
object, had been evicted -- information that is not directly observable
from a single trace replay, because only one candidate is actually evicted
at each decision. This paper studies a reusable way to construct that
counterfactual evidence at scale, and evaluates when it is informative,
when it degenerates, how it relates to actual closed-loop cache
performance, and how sensitive it is to one of its own modeling
assumptions.
\end{sloppypar}

\lafc\ materializes, for real cache traces, a per-decision, per-candidate
finite-horizon counterfactual miss cost: for each object that could have
been evicted at a full-cache miss, how many additional misses would occur
over the next $H$ requests if that object, specifically, had been the one
evicted. This candidate-level quantity is stored as \texttt{y\_loss} (a
convenience alias \texttt{y\_value} $= -\texttt{y\_loss}$ is also provided;
see Section~\ref{sec:tasks-baselines}). Section~\ref{sec:background} works
through one full decision by hand -- resident objects, a forced eviction, a
future request window, and the resulting candidate costs, including a case
where two candidates tie for optimal -- before any aggregate statistic is
introduced, so that every quantity used later in the paper is grounded in
a concrete example first.

\begin{sloppypar}
This paper does not propose a new eviction policy. It evaluates a reusable
offline label against actual closed-loop cache behavior, organized around
five questions. \textbf{(RQ1)} Is the finite-horizon counterfactual target
actually discriminative between candidates, and when does it degenerate to
an uninformative tie? It is not uniformly discriminative: aggregate tie
rates are substantial, and no audited decision has a strict unique winner,
but discriminativeness varies strongly and predictably by workload, cache
capacity, and horizon, and one workload in our evaluation is a genuine,
mechanistically explained negative control. \textbf{(RQ2)} Does the
offline candidate ranking correspond to actual closed-loop cache
performance? Closed-loop replay of seven policies shows offline and
closed-loop orderings agree in the large majority of workload/capacity
regimes where the offline signal itself is not near zero, with two
specific, characterized exceptions. \textbf{(RQ3)} How does this
correspondence depend on workload, capacity, and horizon? Capacity-scale
locality in the underlying trace correlates with both offline
discriminativeness and closed-loop policy separation, and among the three
main-release horizons ($H\in\{4,8,16\}$), $H=16$ aligns offline and
closed-loop evidence most strongly, though modestly; a matched-population
audit extending to $H=128$ on the identical decisions shows informativeness
keeps rising with horizon without eliminating low-information regimes.
\textbf{(RQ4)} How sensitive are the labels to the choice of continuation
policy after a forced eviction? A sampled study and a separate
full-population census both find the canonical LRU-continuation labels
robust to MRU and mean-random continuation, within the capacities and
policies actually tested. \textbf{(RQ5)} What does this imply in practice
for a researcher choosing between offline and closed-loop evaluation? The
offline label complements, rather than replaces, closed-loop evaluation;
Section~\ref{sec:practical-implications} gives concrete guidance for when
it can be trusted.
\end{sloppypar}

Every quantitative claim in this paper traces to a specific, independently
validated evidence artifact, with exact scope, sample sizes, and caveats
stated alongside each result rather than presented as unqualified findings
(Section~\ref{sec:release-validation}). The same underlying labels are
also released as candidate-, decision-, and pairwise-level views at
production scale (Table~\ref{tab:dataset-scale}) for reuse by other
researchers; release construction, schema, and validation are described in
Sections~\ref{sec:benchmark-design}--\ref{sec:release-validation} but are
not this paper's central argument.

To our knowledge, no existing public cache-replacement resource provides
reusable candidate-level counterfactual labels at this scale together with
a direct evaluation of their relationship to closed-loop cache-policy
behavior; this is a specific, bounded claim, not one of exhaustive
prior-art coverage (Section~\ref{sec:related-work}).

The rest of the paper proceeds as follows. Section~\ref{sec:background}
works through the worked example and motivates the counterfactual label.
Sections~\ref{sec:benchmark-design}--\ref{sec:release-validation}
describe construction, label definition, schema, and release validation.
Section~\ref{sec:tasks-baselines} defines the benchmark tasks. Section~\ref{sec:characterization}
reports target discriminativeness (RQ1) and, together with material drawn
from the closed-loop production evaluation, addresses RQ2--RQ4 (offline/
closed-loop correspondence, workload/capacity/horizon dependence, and
continuation-policy robustness); practical implications (RQ5) are
synthesized alongside it. Section~\ref{sec:limitations} states limitations,
ethics, and reuse constraints, followed by related work and conclusion.
